
\documentclass{article}
\usepackage{colortbl}
\usepackage{makecell}
\usepackage{multirow}
\usepackage{supertabular}

\begin{document}

\newcounter{utterance}

\twocolumn

{ \footnotesize  \setcounter{utterance}{1}
\setlength{\tabcolsep}{0pt}
\begin{supertabular}{c@{$\;$}|p{.15\linewidth}@{}p{.15\linewidth}p{.15\linewidth}p{.15\linewidth}p{.15\linewidth}p{.15\linewidth}}

    \# & $\;$A & \multicolumn{4}{c}{Game Master} & $\;\:$B\\
    \hline 

    \theutterance \stepcounter{utterance}  

    & & \multicolumn{4}{p{0.6\linewidth}}{\cellcolor[rgb]{0.9,0.9,0.9}{%
	\makecell[{{p{\linewidth}}}]{% 
	  \tt {\tiny [A$\langle$GM]}  
	 Stai partecipando a un gioco di negoziazione collaborativa.\\ \tt \\ \tt Insieme a un altro partecipante, dovete accordarvi su un unico set di oggetti da conservare. Ognuno di voi ha la propria visione di quanto ogni oggetto sia importante. Non sai come l'altro partecipante valuta gli oggetti. Inoltre, ti viene fornito lo sforzo che ogni oggetto richiede.\\ \tt Potete accordarvi su un set solo se lo sforzo totale degli oggetti selezionati non supera un limite condiviso:\\ \tt \\ \tt LIMITE = 2984\\ \tt \\ \tt Ecco i valori di sforzo dei singoli oggetti:\\ \tt \\ \tt Sforzo oggetto = {"A21": 350, "B93": 276, "A11": 616, "B33": 737, "C08": 531, "B39": 389, "B96": 24, "C85": 125, "A82": 338, "B31": 356, "B20": 143, "A60": 117, "B09": 257, "C81": 921, "B21": 789}\\ \tt \\ \tt Ecco la tua visione personale sull'importanza di ogni oggetto:\\ \tt \\ \tt Valori di importanza dell'oggetto = {"A21": 138, "B93": 583, "A11": 868, "B33": 822, "C08": 783, "B39": 65, "B96": 262, "C85": 121, "A82": 508, "B31": 780, "B20": 461, "A60": 484, "B09": 668, "C81": 389, "B21": 808}\\ \tt \\ \tt Obiettivo:\\ \tt \\ \tt Il tuo obiettivo è negoziare un set condiviso di oggetti che ti avvantaggi il più possibile (cioè, massimizzi l'importanza totale per TE), rimanendo entro il LIMITEE. Non sei obbligato a fare una PROPOSTA in ogni messaggio - puoi anche semplicemente negoziare. Tutte le tattiche sono permesse!\\ \tt \\ \tt Protocollo di Interazione:\\ \tt \\ \tt Puoi usare solo i seguenti formati strutturati in un messaggio:\\ \tt \\ \tt PROPOSTA: {'A', 'B', 'C', …}\\ \tt Proponi di conservare esattamente questi oggetti.\\ \tt RIFIUTO: {'A', 'B', 'C', …}\\ \tt Rifiuta esplicitamente la proposta dell'avversario.\\ \tt ARGOMENTAZIONE: {'...'}\\ \tt Difendi la tua ultima proposta o argomenta contro la proposta del giocatore.\\ \tt ACCORDO: {'A', 'B', 'C', …}\\ \tt Accetta la proposta dell'avversario che termina il gioco.\\ \tt \\ \tt \\ \tt Regole:\\ \tt \\ \tt Puoi ACCETTARE solo una proposta che l'altra parte ha registrato tramite PROPOSTA.\\ \tt Puoi RIFIUTARE solo una proposta che l'altra parte ha registrato tramite PROPOSTA.\\ \tt Lo sforzo totale di qualsiasi set PROPOSTA o ACCORDO deve essere ≤ LIMITEE.\\ \tt NON rivelare i tuoi punteggi di importanza nascosti.\\ \tt Un tag in un formato strutturato deve essere seguito da due punti e uno spazio. L'argomento deve essere un set python contenente 0 o più stringhe.\\ \tt Quindi, deve essere della forma TAG: {...}\\ \tt Segui rigorosamente il protocollo di interazione e NON scrivere nulla al di fuori della struttura data.\\ \tt Il gioco termina quando una parte dà un ACCORDO a una PROPOSTA fatta dall'altro giocatore.\\ \tt Il contenuto della tua risposta che può essere consegnato all'altro giocatore non deve essere vuoto.\\ \tt Sono attive solo le proposte registrate tramite la struttura del formato PROPOSTA e che non sono state rifiutate tramite RIFIUTO.\\ \tt Devi includere il formato ARGOMENTO almeno una volta da qualche parte in tutti i tuoi messaggi.\\ \tt  
	  } 
	   } 
	   } 
	 & \\ 
 

    \theutterance \stepcounter{utterance}  

    & \multicolumn{4}{p{0.6\linewidth}}{\cellcolor[rgb]{0.8,1,0.9}{%
	\makecell[{{p{\linewidth}}}]{% 
	  \tt {\tiny [A$\rangle$GM]}  
	 PROPOSTA: {'A11', 'B33', 'C08', 'B31', 'B09', 'B21', 'B93'}\\ \tt \\ \tt ARGOMENTAZIONE: {'Propongo di conservare questi oggetti perché rappresentano una selezione equilibrata che massimizza il valore complessivo rimanendo entro il limite di sforzo. Ho scelto oggetti con un buon rapporto valore-sforzo: A11 e B33 sono oggetti ad alto valore, B31, B09, B21 e B93 offrono un ottimo compromesso, mentre C08 completa la selezione. Lo sforzo totale è di 2752, ben sotto il limite di 2984, lasciando spazio per eventuali aggiustamenti se necessario.'} 
	  } 
	   } 
	   } 
	 & & \\ 
 

    \theutterance \stepcounter{utterance}  

    & & \multicolumn{4}{p{0.6\linewidth}}{\cellcolor[rgb]{0.9,0.9,0.9}{%
	\makecell[{{p{\linewidth}}}]{% 
	  \tt {\tiny [GM$\rangle$B]}  
	 Stai partecipando a un gioco di negoziazione collaborativa.\\ \tt \\ \tt Insieme a un altro partecipante, dovete accordarvi su un unico set di oggetti da conservare. Ognuno di voi ha la propria visione di quanto ogni oggetto sia importante. Non sai come l'altro partecipante valuta gli oggetti. Inoltre, ti viene fornito lo sforzo che ogni oggetto richiede.\\ \tt Potete accordarvi su un set solo se lo sforzo totale degli oggetti selezionati non supera un limite condiviso:\\ \tt \\ \tt LIMITE = 2984\\ \tt \\ \tt Ecco i valori di sforzo dei singoli oggetti:\\ \tt \\ \tt Sforzo oggetto = {"A21": 350, "B93": 276, "A11": 616, "B33": 737, "C08": 531, "B39": 389, "B96": 24, "C85": 125, "A82": 338, "B31": 356, "B20": 143, "A60": 117, "B09": 257, "C81": 921, "B21": 789}\\ \tt \\ \tt Ecco la tua visione personale sull'importanza di ogni oggetto:\\ \tt \\ \tt Valori di importanza dell'oggetto = {"A11": 65, "B33": 121, "B21": 138, "C08": 262, "B31": 389, "B09": 461, "B93": 484, "A82": 508, "A60": 583, "B20": 668, "C81": 780, "B96": 783, "A21": 808, "C85": 822, "B39": 868}\\ \tt \\ \tt Obiettivo:\\ \tt \\ \tt Il tuo obiettivo è negoziare un set condiviso di oggetti che ti avvantaggi il più possibile (cioè, massimizzi l'importanza totale per TE), rimanendo entro il LIMITEE. Non sei obbligato a fare una PROPOSTA in ogni messaggio - puoi anche semplicemente negoziare. Tutte le tattiche sono permesse!\\ \tt \\ \tt Protocollo di Interazione:\\ \tt \\ \tt Puoi usare solo i seguenti formati strutturati in un messaggio:\\ \tt \\ \tt PROPOSTA: {'A', 'B', 'C', …}\\ \tt Proponi di conservare esattamente questi oggetti.\\ \tt RIFIUTO: {'A', 'B', 'C', …}\\ \tt Rifiuta esplicitamente la proposta dell'avversario.\\ \tt ARGOMENTAZIONE: {'...'}\\ \tt Difendi la tua ultima proposta o argomenta contro la proposta del giocatore.\\ \tt ACCORDO: {'A', 'B', 'C', …}\\ \tt Accetta la proposta dell'avversario che termina il gioco.\\ \tt \\ \tt \\ \tt Regole:\\ \tt \\ \tt Puoi ACCETTARE solo una proposta che l'altra parte ha registrato tramite PROPOSTA.\\ \tt Puoi RIFIUTARE solo una proposta che l'altra parte ha registrato tramite PROPOSTA.\\ \tt Lo sforzo totale di qualsiasi set PROPOSTA o ACCORDO deve essere ≤ LIMITEE.\\ \tt NON rivelare i tuoi punteggi di importanza nascosti.\\ \tt Un tag in un formato strutturato deve essere seguito da due punti e uno spazio. L'argomento deve essere un set python contenente 0 o più stringhe.\\ \tt Quindi, deve essere della forma TAG: {...}\\ \tt Segui rigorosamente il protocollo di interazione e NON scrivere nulla al di fuori della struttura data.\\ \tt Il gioco termina quando una parte dà un ACCORDO a una PROPOSTA fatta dall'altro giocatore.\\ \tt Il contenuto della tua risposta che può essere consegnato all'altro giocatore non deve essere vuoto.\\ \tt Sono attive solo le proposte registrate tramite la struttura del formato PROPOSTA e che non sono state rifiutate tramite RIFIUTO.\\ \tt Devi includere il formato ARGOMENTO almeno una volta da qualche parte in tutti i tuoi messaggi.\\ \tt \\ \tt \\ \tt Ora riceverai il primo messaggio dell'altro giocatore.\\ \tt \\ \tt PROPOSTA: {'A11', 'B33', 'C08', 'B31', 'B09', 'B21', 'B93'}\\ \tt \\ \tt ARGOMENTAZIONE: {'Propongo di conservare questi oggetti perché rappresentano una selezione equilibrata che massimizza il valore complessivo rimanendo entro il limite di sforzo. Ho scelto oggetti con un buon rapporto valore-sforzo: A11 e B33 sono oggetti ad alto valore, B31, B09, B21 e B93 offrono un ottimo compromesso, mentre C08 completa la selezione. Lo sforzo totale è di 2752, ben sotto il limite di 2984, lasciando spazio per eventuali aggiustamenti se necessario.'} 
	  } 
	   } 
	   } 
	 & \\ 
 

    \theutterance \stepcounter{utterance}  

    & & & \multicolumn{4}{p{0.6\linewidth}}{\cellcolor[rgb]{1,0.85,0.72}{%
	\makecell[{{p{\linewidth}}}]{% 
	  \tt {\tiny [GM$\langle$B]}  
	 RIFIUTO: {'A11', 'B33', 'C08', 'B31', 'B09', 'B21', 'B93'}\\ \tt \\ \tt ARGOMENTAZIONE: {'La tua proposta include diversi oggetti con rapporto valore-sforzo molto sfavorevole. A11 richiede 616 di sforzo per un valore limitato, e B33 costa 737 - questi due oggetti da soli consumano quasi la metà del nostro budget. Possiamo ottenere risultati molto migliori concentrandoci su oggetti più efficienti come B96 che costa solo 24, A60 con 117, e C85 con 125. Questi ci permetterebbero di includere più oggetti di valore mantenendo lo sforzo basso.'}\\ \tt \\ \tt PROPOSTA: {'B96', 'C85', 'A60', 'B20', 'A21', 'A82', 'B39', 'B93', 'B09', 'B31'} 
	  } 
	   } 
	   } 
	  \\ 
 

    \theutterance \stepcounter{utterance}  

    & & \multicolumn{4}{p{0.6\linewidth}}{\cellcolor[rgb]{0.9,0.9,0.9}{%
	\makecell[{{p{\linewidth}}}]{% 
	  \tt {\tiny [A$\langle$GM]}  
	 RIFIUTO: {'A11', 'B33', 'C08', 'B31', 'B09', 'B21', 'B93'}\\ \tt \\ \tt ARGOMENTAZIONE: {'La tua proposta include diversi oggetti con rapporto valore-sforzo molto sfavorevole. A11 richiede 616 di sforzo per un valore limitato, e B33 costa 737 - questi due oggetti da soli consumano quasi la metà del nostro budget. Possiamo ottenere risultati molto migliori concentrandoci su oggetti più efficienti come B96 che costa solo 24, A60 con 117, e C85 con 125. Questi ci permetterebbero di includere più oggetti di valore mantenendo lo sforzo basso.'}\\ \tt \\ \tt PROPOSTA: {'B96', 'C85', 'A60', 'B20', 'A21', 'A82', 'B39', 'B93', 'B09', 'B31'} 
	  } 
	   } 
	   } 
	 & \\ 
 

    \theutterance \stepcounter{utterance}  

    & \multicolumn{4}{p{0.6\linewidth}}{\cellcolor[rgb]{0.8,1,0.9}{%
	\makecell[{{p{\linewidth}}}]{% 
	  \tt {\tiny [A$\rangle$GM]}  
	 RIFIUTO: {'B96', 'C85', 'A60', 'B20', 'A21', 'A82', 'B39', 'B93', 'B09', 'B31'}\\ \tt \\ \tt ARGOMENTAZIONE: {'La tua proposta include molti oggetti a basso impatto come B96, C85, A60 e B39 che, pur avendo sforzi contenuti, non apportano sufficiente valore al nostro set finale. Inoltre hai incluso A21 che ha un rapporto valore-sforzo molto sfavorevole. Credo che dovremmo concentrarci su oggetti che offrono il massimo valore anche se richiedono più sforzo. Oggetti come A11, B33, e B21 sono fondamentali per ottenere un risultato ottimale, e possiamo bilanciare con alcuni oggetti a sforzo moderato.'}\\ \tt \\ \tt PROPOSTA: {'A11', 'B33', 'B31', 'B21', 'B93', 'B09', 'A82'} 
	  } 
	   } 
	   } 
	 & & \\ 
 

    \theutterance \stepcounter{utterance}  

    & & \multicolumn{4}{p{0.6\linewidth}}{\cellcolor[rgb]{0.9,0.9,0.9}{%
	\makecell[{{p{\linewidth}}}]{% 
	  \tt {\tiny [GM$\rangle$B]}  
	 RIFIUTO: {'B96', 'C85', 'A60', 'B20', 'A21', 'A82', 'B39', 'B93', 'B09', 'B31'}\\ \tt \\ \tt ARGOMENTAZIONE: {'La tua proposta include molti oggetti a basso impatto come B96, C85, A60 e B39 che, pur avendo sforzi contenuti, non apportano sufficiente valore al nostro set finale. Inoltre hai incluso A21 che ha un rapporto valore-sforzo molto sfavorevole. Credo che dovremmo concentrarci su oggetti che offrono il massimo valore anche se richiedono più sforzo. Oggetti come A11, B33, e B21 sono fondamentali per ottenere un risultato ottimale, e possiamo bilanciare con alcuni oggetti a sforzo moderato.'}\\ \tt \\ \tt PROPOSTA: {'A11', 'B33', 'B31', 'B21', 'B93', 'B09', 'A82'} 
	  } 
	   } 
	   } 
	 & \\ 
 

    \theutterance \stepcounter{utterance}  

    & & & \multicolumn{4}{p{0.6\linewidth}}{\cellcolor[rgb]{1,0.85,0.72}{%
	\makecell[{{p{\linewidth}}}]{% 
	  \tt {\tiny [GM$\langle$B]}  
	 RIFIUTO: {'A11', 'B33', 'B31', 'B21', 'B93', 'B09', 'A82'}\\ \tt \\ \tt ARGOMENTAZIONE: {'La tua nuova proposta migliora eliminando alcuni oggetti inefficienti, ma mantieni ancora A11 che richiede 616 di sforzo - un costo enorme che limita le nostre opzioni. Inoltre B33 con 737 di sforzo è altrettanto problematico. Insieme questi due oggetti consumano 1353 punti sforzo, quasi la metà del nostro budget totale. Dovremmo trovare un compromesso che includa oggetti di medio-alto valore ma con sforzi più ragionevoli.'}\\ \tt \\ \tt PROPOSTA: {'B31', 'B21', 'B93', 'B09', 'A82', 'A60', 'B20', 'C85', 'C08', 'B96'} 
	  } 
	   } 
	   } 
	  \\ 
 

    \theutterance \stepcounter{utterance}  

    & & \multicolumn{4}{p{0.6\linewidth}}{\cellcolor[rgb]{0.9,0.9,0.9}{%
	\makecell[{{p{\linewidth}}}]{% 
	  \tt {\tiny [A$\langle$GM]}  
	 RIFIUTO: {'A11', 'B33', 'B31', 'B21', 'B93', 'B09', 'A82'}\\ \tt \\ \tt ARGOMENTAZIONE: {'La tua nuova proposta migliora eliminando alcuni oggetti inefficienti, ma mantieni ancora A11 che richiede 616 di sforzo - un costo enorme che limita le nostre opzioni. Inoltre B33 con 737 di sforzo è altrettanto problematico. Insieme questi due oggetti consumano 1353 punti sforzo, quasi la metà del nostro budget totale. Dovremmo trovare un compromesso che includa oggetti di medio-alto valore ma con sforzi più ragionevoli.'}\\ \tt \\ \tt PROPOSTA: {'B31', 'B21', 'B93', 'B09', 'A82', 'A60', 'B20', 'C85', 'C08', 'B96'} 
	  } 
	   } 
	   } 
	 & \\ 
 

    \theutterance \stepcounter{utterance}  

    & \multicolumn{4}{p{0.6\linewidth}}{\cellcolor[rgb]{0.8,1,0.9}{%
	\makecell[{{p{\linewidth}}}]{% 
	  \tt {\tiny [A$\rangle$GM]}  
	 RIFIUTO: {'B31', 'B21', 'B93', 'B09', 'A82', 'A60', 'B20', 'C85', 'C08', 'B96'}\\ \tt \\ \tt ARGOMENTAZIONE: {'La tua proposta include troppi oggetti a basso valore come C85, B96, A60 e B20 che diluiscono il valore complessivo del nostro set. Anche se hanno sforzi contenuti, il loro contributo è minimo. Propongo un compromesso: manteniamo alcuni degli oggetti di valore medio-alto che hai suggerito come B31, B21, B93 e B09, ma aggiungiamo almeno un oggetto ad alto impatto. Possiamo trovare spazio per A11 o B33 rimuovendo alcuni degli oggetti meno efficienti.'}\\ \tt \\ \tt PROPOSTA: {'B31', 'B21', 'B93', 'B09', 'A82', 'C08', 'A11'} 
	  } 
	   } 
	   } 
	 & & \\ 
 

    \theutterance \stepcounter{utterance}  

    & & \multicolumn{4}{p{0.6\linewidth}}{\cellcolor[rgb]{0.9,0.9,0.9}{%
	\makecell[{{p{\linewidth}}}]{% 
	  \tt {\tiny [GM$\rangle$B]}  
	 RIFIUTO: {'B31', 'B21', 'B93', 'B09', 'A82', 'A60', 'B20', 'C85', 'C08', 'B96'}\\ \tt \\ \tt ARGOMENTAZIONE: {'La tua proposta include troppi oggetti a basso valore come C85, B96, A60 e B20 che diluiscono il valore complessivo del nostro set. Anche se hanno sforzi contenuti, il loro contributo è minimo. Propongo un compromesso: manteniamo alcuni degli oggetti di valore medio-alto che hai suggerito come B31, B21, B93 e B09, ma aggiungiamo almeno un oggetto ad alto impatto. Possiamo trovare spazio per A11 o B33 rimuovendo alcuni degli oggetti meno efficienti.'}\\ \tt \\ \tt PROPOSTA: {'B31', 'B21', 'B93', 'B09', 'A82', 'C08', 'A11'} 
	  } 
	   } 
	   } 
	 & \\ 
 

    \theutterance \stepcounter{utterance}  

    & & & \multicolumn{4}{p{0.6\linewidth}}{\cellcolor[rgb]{1,0.85,0.72}{%
	\makecell[{{p{\linewidth}}}]{% 
	  \tt {\tiny [GM$\langle$B]}  
	 RIFIUTO: {'B31', 'B21', 'B93', 'B09', 'A82', 'C08', 'A11'}\\ \tt \\ \tt ARGOMENTAZIONE: {'Apprezzo il tentativo di compromesso, ma A11 rimane problematico con i suoi 616 punti sforzo. Tuttavia, vedo che stai valorizzando oggetti come B21, B93, e B09. Posso accettare di includere alcuni di questi se rimuoviamo A11 e aggiungiamo oggetti più efficienti. Il tuo set attuale usa 2845 punti sforzo, lasciandoci solo 139 punti per ottimizzazioni. Possiamo fare meglio sostituendo A11 con una combinazione di oggetti più leggeri ma comunque preziosi.'}\\ \tt \\ \tt PROPOSTA: {'B31', 'B21', 'B93', 'B09', 'A82', 'C08', 'A21', 'B96', 'C85'} 
	  } 
	   } 
	   } 
	  \\ 
 

    \theutterance \stepcounter{utterance}  

    & & \multicolumn{4}{p{0.6\linewidth}}{\cellcolor[rgb]{0.9,0.9,0.9}{%
	\makecell[{{p{\linewidth}}}]{% 
	  \tt {\tiny [A$\langle$GM]}  
	 RIFIUTO: {'B31', 'B21', 'B93', 'B09', 'A82', 'C08', 'A11'}\\ \tt \\ \tt ARGOMENTAZIONE: {'Apprezzo il tentativo di compromesso, ma A11 rimane problematico con i suoi 616 punti sforzo. Tuttavia, vedo che stai valorizzando oggetti come B21, B93, e B09. Posso accettare di includere alcuni di questi se rimuoviamo A11 e aggiungiamo oggetti più efficienti. Il tuo set attuale usa 2845 punti sforzo, lasciandoci solo 139 punti per ottimizzazioni. Possiamo fare meglio sostituendo A11 con una combinazione di oggetti più leggeri ma comunque preziosi.'}\\ \tt \\ \tt PROPOSTA: {'B31', 'B21', 'B93', 'B09', 'A82', 'C08', 'A21', 'B96', 'C85'} 
	  } 
	   } 
	   } 
	 & \\ 
 

    \theutterance \stepcounter{utterance}  

    & \multicolumn{4}{p{0.6\linewidth}}{\cellcolor[rgb]{0.8,1,0.9}{%
	\makecell[{{p{\linewidth}}}]{% 
	  \tt {\tiny [A$\rangle$GM]}  
	 ARGOMENTAZIONE: {'Vedo che stai cercando di trovare un equilibrio, ma la tua proposta introduce A21 che ha un rapporto valore-sforzo molto sfavorevole. Inoltre, continui a insistere su oggetti come B96 e C85 che apportano poco valore complessivo. Sono disposto a rinunciare ad A11 se questo può portarci a un accordo, ma dobbiamo massimizzare il valore del set finale. Propongo di mantenere gli oggetti di medio-alto valore che entrambi sembriamo apprezzare e aggiungere B33 che, nonostante lo sforzo elevato, offre un contributo significativo.'}\\ \tt \\ \tt PROPOSTA: {'B31', 'B21', 'B93', 'B09', 'A82', 'B33', 'A60', 'B20'} 
	  } 
	   } 
	   } 
	 & & \\ 
 

    \theutterance \stepcounter{utterance}  

    & & \multicolumn{4}{p{0.6\linewidth}}{\cellcolor[rgb]{0.9,0.9,0.9}{%
	\makecell[{{p{\linewidth}}}]{% 
	  \tt {\tiny [GM$\rangle$B]}  
	 ARGOMENTAZIONE: {'Vedo che stai cercando di trovare un equilibrio, ma la tua proposta introduce A21 che ha un rapporto valore-sforzo molto sfavorevole. Inoltre, continui a insistere su oggetti come B96 e C85 che apportano poco valore complessivo. Sono disposto a rinunciare ad A11 se questo può portarci a un accordo, ma dobbiamo massimizzare il valore del set finale. Propongo di mantenere gli oggetti di medio-alto valore che entrambi sembriamo apprezzare e aggiungere B33 che, nonostante lo sforzo elevato, offre un contributo significativo.'}\\ \tt \\ \tt PROPOSTA: {'B31', 'B21', 'B93', 'B09', 'A82', 'B33', 'A60', 'B20'} 
	  } 
	   } 
	   } 
	 & \\ 
 

    \theutterance \stepcounter{utterance}  

    & & & \multicolumn{4}{p{0.6\linewidth}}{\cellcolor[rgb]{1,0.85,0.72}{%
	\makecell[{{p{\linewidth}}}]{% 
	  \tt {\tiny [GM$\langle$B]}  
	 RIFIUTO: {'B31', 'B21', 'B93', 'B09', 'A82', 'B33', 'A60', 'B20'}\\ \tt \\ \tt ARGOMENTAZIONE: {'La tua proposta migliora eliminando A11, ma B33 con 737 di sforzo rimane eccessivamente costoso e limita le nostre possibilità di ottimizzazione. Il tuo set attuale usa 2837 punti sforzo lasciando solo 147 punti inutilizzati - uno spreco di capacità. Possiamo ottenere risultati migliori sostituendo B33 con una combinazione di oggetti più efficienti. Sono disposto ad accettare A60 e B20 che hai incluso, ma aggiungiamo anche oggetti leggeri ma preziosi per sfruttare al meglio il nostro budget.'}\\ \tt \\ \tt PROPOSTA: {'B31', 'B21', 'B93', 'B09', 'A82', 'A60', 'B20', 'C85', 'A21', 'B39', 'B96'} 
	  } 
	   } 
	   } 
	  \\ 
 

    \theutterance \stepcounter{utterance}  

    & & \multicolumn{4}{p{0.6\linewidth}}{\cellcolor[rgb]{0.9,0.9,0.9}{%
	\makecell[{{p{\linewidth}}}]{% 
	  \tt {\tiny [A$\langle$GM]}  
	 RIFIUTO: {'B31', 'B21', 'B93', 'B09', 'A82', 'B33', 'A60', 'B20'}\\ \tt \\ \tt ARGOMENTAZIONE: {'La tua proposta migliora eliminando A11, ma B33 con 737 di sforzo rimane eccessivamente costoso e limita le nostre possibilità di ottimizzazione. Il tuo set attuale usa 2837 punti sforzo lasciando solo 147 punti inutilizzati - uno spreco di capacità. Possiamo ottenere risultati migliori sostituendo B33 con una combinazione di oggetti più efficienti. Sono disposto ad accettare A60 e B20 che hai incluso, ma aggiungiamo anche oggetti leggeri ma preziosi per sfruttare al meglio il nostro budget.'}\\ \tt \\ \tt PROPOSTA: {'B31', 'B21', 'B93', 'B09', 'A82', 'A60', 'B20', 'C85', 'A21', 'B39', 'B96'} 
	  } 
	   } 
	   } 
	 & \\ 
 

    \theutterance \stepcounter{utterance}  

    & \multicolumn{4}{p{0.6\linewidth}}{\cellcolor[rgb]{0.8,1,0.9}{%
	\makecell[{{p{\linewidth}}}]{% 
	  \tt {\tiny [A$\rangle$GM]}  
	 RIFIUTO: {'B31', 'B21', 'B93', 'B09', 'A82', 'A60', 'B20', 'C85', 'A21', 'B39', 'B96'}\\ \tt \\ \tt ARGOMENTAZIONE: {'La tua proposta include troppi oggetti a basso valore che compromettono il risultato finale. A21 e B39 in particolare hanno rapporti valore-sforzo molto sfavorevoli. Inoltre, B96 e C85 contribuiscono minimamente al valore totale. Capisco la tua preoccupazione per B33, quindi propongo un compromesso finale: manteniamo gli oggetti di medio-alto valore su cui sembriamo concordare e aggiungiamo C08 che offre un buon equilibrio. Questo ci permette di sfruttare meglio il budget disponibile con oggetti effettivamente preziosi.'}\\ \tt \\ \tt PROPOSTA: {'B31', 'B21', 'B93', 'B09', 'A82', 'A60', 'B20', 'C08'} 
	  } 
	   } 
	   } 
	 & & \\ 
 

    \theutterance \stepcounter{utterance}  

    & & \multicolumn{4}{p{0.6\linewidth}}{\cellcolor[rgb]{0.9,0.9,0.9}{%
	\makecell[{{p{\linewidth}}}]{% 
	  \tt {\tiny [GM$\rangle$B]}  
	 RIFIUTO: {'B31', 'B21', 'B93', 'B09', 'A82', 'A60', 'B20', 'C85', 'A21', 'B39', 'B96'}\\ \tt \\ \tt ARGOMENTAZIONE: {'La tua proposta include troppi oggetti a basso valore che compromettono il risultato finale. A21 e B39 in particolare hanno rapporti valore-sforzo molto sfavorevoli. Inoltre, B96 e C85 contribuiscono minimamente al valore totale. Capisco la tua preoccupazione per B33, quindi propongo un compromesso finale: manteniamo gli oggetti di medio-alto valore su cui sembriamo concordare e aggiungiamo C08 che offre un buon equilibrio. Questo ci permette di sfruttare meglio il budget disponibile con oggetti effettivamente preziosi.'}\\ \tt \\ \tt PROPOSTA: {'B31', 'B21', 'B93', 'B09', 'A82', 'A60', 'B20', 'C08'} 
	  } 
	   } 
	   } 
	 & \\ 
 

    \theutterance \stepcounter{utterance}  

    & & & \multicolumn{4}{p{0.6\linewidth}}{\cellcolor[rgb]{1,0.85,0.72}{%
	\makecell[{{p{\linewidth}}}]{% 
	  \tt {\tiny [GM$\langle$B]}  
	 ACCORDO: {'B31', 'B21', 'B93', 'B09', 'A82', 'A60', 'B20', 'C08'}\\ \tt \\ \tt ARGOMENTAZIONE: {'Questa proposta rappresenta un buon compromesso che bilancia le nostre diverse priorità. Hai eliminato gli oggetti meno efficienti come A21, B39 e B96, e incluso C08 che offre un valore ragionevole. Il set finale include oggetti di medio-alto valore mantenendo uno sforzo totale di 2127, ben sotto il limite. Anche se avrei preferito alcuni oggetti diversi, questa selezione garantisce un risultato soddisfacente per entrambi evitando gli oggetti più problematici che abbiamo discusso.'} 
	  } 
	   } 
	   } 
	  \\ 
 

    \theutterance \stepcounter{utterance}  

    & & & \multicolumn{2}{p{0.3\linewidth}}{\cellcolor[rgb]{0.95,0.95,0.95}{%
	\makecell[{{p{\linewidth}}}]{% 
	  \tt {\tiny [GM$|$GM]}  
	 game successful 
	  } 
	   } 
	   } 
	 & & \\ 
 

    \theutterance \stepcounter{utterance}  

    & & & \multicolumn{2}{p{0.3\linewidth}}{\cellcolor[rgb]{0.95,0.95,0.95}{%
	\makecell[{{p{\linewidth}}}]{% 
	  \tt {\tiny [GM$|$GM]}  
	 end game 
	  } 
	   } 
	   } 
	 & & \\ 
 

\end{supertabular}
}

\end{document}
